\documentclass[aspectratio=169]{beamer}

% Tema UFS -- Universidade Federal de Sergipe
\usetheme{UFS}

\usepackage[utf8]{inputenc}
\usepackage[T1]{fontenc}
\usepackage[portuguese]{babel}
\usepackage{amsmath,amssymb}
\usepackage{graphicx}
\usepackage{booktabs}
\usepackage{xcolor}

\usepackage{tikz}
\usetikzlibrary{shapes.geometric, arrows.meta, positioning, matrix, fit,
  backgrounds, decorations.pathreplacing, calc}

% ----- Listagens: paleta VS Code Light+ -----
\usepackage{listings}
\definecolor{bgcolor}{HTML}{FFFFFF}
\definecolor{linecolor}{HTML}{DDDDDD}
\definecolor{numcolor}{HTML}{999999}
\definecolor{kwcolor}{HTML}{0000FF}
\definecolor{strcolor}{HTML}{A31515}
\definecolor{cmtcolor}{HTML}{008000}

\lstdefinestyle{ufsStyle}{
  language=Python,
  backgroundcolor=\color{bgcolor},
  basicstyle=\ttfamily\scriptsize\color{black},
  keywordstyle=\color{kwcolor}\bfseries,
  stringstyle=\color{strcolor},
  commentstyle=\color{cmtcolor}\itshape,
  numberstyle=\tiny\color{numcolor},
  numbers=left, numbersep=8pt,
  xleftmargin=16pt, framexleftmargin=16pt,
  breaklines=true, tabsize=4,
  keepspaces=true, showspaces=false,
  showstringspaces=false, showtabs=false,
  frame=single, rulecolor=\color{linecolor},
  framesep=3pt, captionpos=b, upquote=true,
  columns=fullflexible,
  extendedchars=true,
  literate=%
    {á}{{\'a}}1 {à}{{\`a}}1 {ã}{{\~a}}1 {â}{{\^a}}1
    {é}{{\'e}}1 {ê}{{\^e}}1 {í}{{\'i}}1
    {ó}{{\'o}}1 {ô}{{\^o}}1 {õ}{{\~o}}1
    {ú}{{\'u}}1 {ç}{{\c c}}1
    {Á}{{\'A}}1 {É}{{\'E}}1 {Í}{{\'I}}1 {Ó}{{\'O}}1 {Ú}{{\'U}}1 {Ç}{{\c C}}1
}
\lstset{style=ufsStyle}

% Estilo enxuto para saidas de terminal
\lstdefinestyle{saida}{
  language=,
  basicstyle=\ttfamily\scriptsize\color{black},
  numbers=none,
  frame=single, rulecolor=\color{linecolor},
  backgroundcolor=\color{UFSLightGray},
  xleftmargin=0pt, framexleftmargin=0pt,
  breaklines=true,
}

\newcommand{\code}[1]{\texttt{\small #1}}

\title{Condicionais em Python}
\subtitle{Decidindo com \texttt{if}, \texttt{elif} e \texttt{else}}
\author{Prof.\ André Yoshiaki Kashiwabara}
\institute{DCOMP -- Universidade Federal de Sergipe\\
\texttt{andreyoshiaki@dcomp.ufs.br}}
\date{}

\begin{document}
\shorthandoff{"}

\begin{frame}[plain]
  \titlepage
\end{frame}

\begin{frame}{Agenda}
  \tableofcontents
\end{frame}

% ===================================================================
\section{O problema da triagem}
% ===================================================================

\begin{frame}{Uma manhã no pronto-socorro}
  \begin{columns}[T]
    \begin{column}{0.52\textwidth}
      \textbf{A situação}\\[0.2cm]
      Chegam dezenas de pacientes por hora. Alguém precisa decidir,
      para \emph{cada um}, quem é atendido primeiro.

      \vspace{0.3cm}
      A decisão não pode depender do humor de quem está na recepção:
      dois pacientes com os mesmos dados têm de receber a mesma
      classificação.
    \end{column}
    \begin{column}{0.44\textwidth}
      \begin{block}{O que o computador precisa}
        Uma regra escrita de forma que a máquina consiga
        \textbf{seguir}: dados de entrada, testes, uma resposta.
      \end{block}
      \begin{exampleblock}{Nossa ferramenta de hoje}
        O comando \code{if} --- e os seus dois companheiros.
      \end{exampleblock}
    \end{column}
  \end{columns}
\end{frame}

\begin{frame}{A regra da triagem}
  Cada paciente tem dois dados: \textbf{pressão sistólica}
  (real, na escala mmHg/10 --- o valor \code{14.0} representa 140~mmHg)
  e \textbf{idade} (inteiro, em anos).

  \vspace{0.25cm}
  \begin{center}
  \scriptsize
  \begin{tabular}{@{}ll@{}}
    \toprule
    \textbf{Classificação} & \textbf{Descrição} \\
    \midrule
    \code{NORMAL}  & pressão e idade dentro dos limites saudáveis \\
    \code{BAIXO}   & pressão acima de 120 mmHg, sem outros agravantes \\
    \code{MEDIO}   & pressão acima de 140 mmHg \textbf{ou} idade acima de 60 \\
    \code{ALTO}    & pressão acima de 140 mmHg \textbf{e} idade acima de 60 \\
    \code{CRITICO} & pressão inferior a 90 mmHg --- emergência, qualquer idade \\
    \bottomrule
  \end{tabular}
  \end{center}

  \vspace{0.15cm}
  \begin{alertblock}{O aviso do enunciado}
    ``Um mesmo paciente pode se enquadrar em mais de uma descrição;
    cabe a você determinar a \textbf{ordem correta de verificação}.''
  \end{alertblock}
\end{frame}

\begin{frame}{Os cinco pacientes desta manhã}
  \begin{center}
  \begin{tabular}{@{}lrrl@{}}
    \toprule
    \textbf{Nome} & \textbf{Pressão} & \textbf{Idade} & \textbf{Em mmHg} \\
    \midrule
    Antonio & 15.0 & 70 & 150 mmHg \\
    Beatriz & 14.5 & 30 & 145 mmHg \\
    Cesar   & 12.5 & 80 & 125 mmHg \\
    Daniela & 11.0 & 20 & 110 mmHg \\
    Eduardo &  8.5 & 45 &  85 mmHg \\
    \bottomrule
  \end{tabular}
  \end{center}

  \vspace{0.35cm}
  \begin{block}{Faça agora, no papel}
    Classifique os cinco pacientes usando a tabela do slide anterior.
    Guarde a sua folha --- vamos comparar com o computador.
  \end{block}
\end{frame}

\begin{frame}{O paciente que complica tudo}
  \textbf{Antonio: 150 mmHg, 70 anos.} Confira as descrições:

  \vspace{0.25cm}
  \begin{center}
  \small
  \begin{tabular}{@{}lll@{}}
    \toprule
    \textbf{Descrição} & \textbf{Serve para o Antonio?} & \textbf{Por quê} \\
    \midrule
    \code{BAIXO}   & sim & 150 está acima de 120 \\
    \code{MEDIO}   & sim & 150 está acima de 140 \\
    \code{ALTO}    & sim & 150 acima de 140 \textbf{e} 70 acima de 60 \\
    \bottomrule
  \end{tabular}
  \end{center}

  \vspace{0.25cm}
  \begin{alertblock}{A pergunta central da aula}
    Três descrições verdadeiras ao mesmo tempo. O computador vai
    escolher \textbf{uma}. Quem decide qual? \\
    Resposta: quem escreveu o código --- pela \textbf{ordem} em que
    colocou os testes.
  \end{alertblock}
\end{frame}

\begin{frame}{O mapa das decisões}
  \begin{center}
  \begin{tikzpicture}[x=0.092cm, y=0.038cm, font=\tiny]
    % regioes
    \fill[UFSRed!70]     (80,0)  rectangle (90,90);
    \fill[green!35!white](90,0)  rectangle (120,60);
    \fill[UFSYellow!60]  (120,0) rectangle (140,60);
    \fill[orange!55]     (140,0) rectangle (170,60);
    \fill[orange!55]     (90,60) rectangle (140,90);
    \fill[UFSRed!40]     (140,60)rectangle (170,90);

    % grade
    \draw[gray!60] (80,0) rectangle (170,90);
    \draw[gray!60] (90,0)--(90,90);
    \draw[gray!60] (120,0)--(120,60);
    \draw[gray!60] (140,0)--(140,90);
    \draw[gray!60] (90,60)--(170,60);

    % rotulos das regioes
    \node at (85,72)   {\rotatebox{90}{\textbf{CRITICO}}};
    \node at (100,45)  {\textbf{NORMAL}};
    \node at (130,45)  {\textbf{BAIXO}};
    \node at (160,14)  {\textbf{MEDIO}};
    \node at (103,70)  {\textbf{MEDIO}};
    \node at (163,83)  {\textbf{ALTO}};

    % eixos
    \draw[-{Latex[length=2mm]},thick] (80,0) -- (176,0)
      node[right,font=\scriptsize] {pressão (mmHg)};
    \draw[-{Latex[length=2mm]},thick] (80,0) -- (80,98)
      node[above,font=\scriptsize] {idade};
    \foreach \x in {90,120,140} \node[below] at (\x,0) {\x};
    \node[left] at (80,60) {60};

    % pacientes
    \foreach \px/\py/\nm in {150/70/Antonio, 145/30/Beatriz, 125/80/Cesar,
                             110/20/Daniela, 85/45/Eduardo} {
      \fill (\px,\py) circle (1.6pt);
    }
    \node[below right, font=\tiny] at (150,70) {Antonio};
    \node[above right, font=\tiny] at (145,30) {Beatriz};
    \node[above,       font=\tiny] at (125,80) {Cesar};
    \node[above right, font=\tiny] at (110,20) {Daniela};
    \node[above right, font=\tiny] at (85,45)  {Eduardo};
  \end{tikzpicture}
  \end{center}

  \vspace{-0.2cm}
  {\small Cada ponto é um paciente. Cada região, uma classificação.
  Repare: o Antonio está no canto onde \emph{pressão alta} e
  \emph{idade alta} se encontram.}
\end{frame}

\begin{frame}[fragile,shrink=8]{Antes de rodar: qual será a saída?}
  \begin{columns}[T]
    \begin{column}{0.6\textwidth}
\begin{lstlisting}
def risco(pressao, idade):
    if pressao < 9.0:
        return "CRITICO"
    elif pressao > 14.0 and idade > 60:
        return "ALTO"
    elif pressao > 14.0 or idade > 60:
        return "MEDIO"
    elif pressao > 12.0:
        return "BAIXO"
    else:
        return "NORMAL"

print("Antonio:", risco(15.0, 70))
print("Beatriz:", risco(14.5, 30))
print("Cesar:",   risco(12.5, 80))
print("Daniela:", risco(11.0, 20))
print("Eduardo:", risco(8.5,  45))
\end{lstlisting}
    \end{column}
    \begin{column}{0.36\textwidth}
      \begin{block}{Escreva a sua previsão}
        \code{Antonio: \rule{1.8cm}{0.4pt}}\\[0.25cm]
        \code{Beatriz: \rule{1.8cm}{0.4pt}}\\[0.25cm]
        \code{Cesar: \rule{1.8cm}{0.4pt}}\\[0.25cm]
        \code{Daniela: \rule{1.8cm}{0.4pt}}\\[0.25cm]
        \code{Eduardo: \rule{1.8cm}{0.4pt}}
      \end{block}
      \begin{exampleblock}{Regra do jogo}
        {\small Escreva \textbf{antes} de rodar. Errar a previsão
        é o momento em que se aprende mais.}
      \end{exampleblock}
    \end{column}
  \end{columns}
\end{frame}

% ===================================================================
\section{Como o computador decide}
% ===================================================================

\begin{frame}[fragile]{A saída real}
  \begin{columns}[T]
    \begin{column}{0.45\textwidth}
      \begin{block}{O que o Python imprimiu}
\begin{lstlisting}[style=saida]
Antonio: ALTO
Beatriz: MEDIO
Cesar: MEDIO
Daniela: NORMAL
Eduardo: CRITICO
\end{lstlisting}
      \end{block}
    \end{column}
    \begin{column}{0.5\textwidth}
      \begin{exampleblock}{Compare com a sua folha}
        \begin{itemize}
          \item Acertou os cinco? Ótimo --- agora explique \emph{por que}
                o Antonio é \code{ALTO} e não \code{MEDIO}.
          \item Errou algum? Marque qual. Vamos descobrir exatamente
                em que linha o computador decidiu.
        \end{itemize}
      \end{exampleblock}
      \begin{block}{Onde mora a decisão}
        A resposta está nas quatro palavras \code{if}, \code{elif},
        \code{or}, \code{and} --- e na ordem delas.
      \end{block}
    \end{column}
  \end{columns}
\end{frame}

\begin{frame}[shrink=5]{Anatomia de um \texttt{if}}
  \begin{center}
  \begin{tikzpicture}[font=\small]
    \node[font=\ttfamily\large] (l1) at (0,0) {if pressao < 9.0:};
    \node[font=\ttfamily\large, anchor=west] (l2) at (0.9,-1.1)
      {return "CRITICO"};

    \draw[UFSBlue, thick, -{Latex[length=2mm]}] (-2.2,1.3) -- (-2.2,0.35);
    \node[UFSBlue, above, font=\scriptsize] at (-2.2,1.3) {palavra-chave};

    \draw[UFSRed, thick, -{Latex[length=2mm]}] (0.1,1.3) -- (0.1,0.35);
    \node[UFSRed, above, font=\scriptsize] at (0.1,1.3) {condição (dá True ou False)};

    \draw[orange!80!black, thick, -{Latex[length=2mm]}] (2.5,1.3) -- (2.45,0.35);
    \node[orange!80!black, above, font=\scriptsize] at (2.7,1.3) {dois-pontos};

    \draw[UFSBlue, thick, -{Latex[length=2mm]}] (5.6,-1.1) -- (3.7,-1.1);
    \node[UFSBlue, right, font=\scriptsize] at (5.6,-1.1)
      {bloco: só roda se a condição for True};

    \draw[decorate, decoration={brace, amplitude=5pt}, UFSGray, thick]
      (0.85,-1.45) -- (0.05,-1.45);
    \node[UFSGray, below, font=\scriptsize, align=center] at (0.45,-1.7)
      {indentação\\(4 espaços)};
  \end{tikzpicture}
  \end{center}

  \vspace{0.2cm}
  \begin{alertblock}{Em Python a indentação não é enfeite}
    Ela é que diz o que está dentro do \code{if}. Perder um espaço
    muda o programa.
  \end{alertblock}
\end{frame}

\begin{frame}[fragile]{A condição é uma pergunta de sim ou não}
  Toda condição é uma expressão que vale \code{True} ou \code{False}.

  \vspace{0.3cm}
  \begin{columns}[T]
    \begin{column}{0.46\textwidth}
      \begin{center}
      \small
      \begin{tabular}{@{}ll@{}}
        \toprule
        \textbf{Operador} & \textbf{Lê-se} \\
        \midrule
        \code{<}  & menor que \\
        \code{<=} & menor ou igual a \\
        \code{>}  & maior que \\
        \code{>=} & maior ou igual a \\
        \code{==} & é igual a \\
        \code{!=} & é diferente de \\
        \bottomrule
      \end{tabular}
      \end{center}
    \end{column}
    \begin{column}{0.5\textwidth}
\begin{lstlisting}
>>> pressao = 15.0
>>> pressao < 9.0
False
>>> pressao > 14.0
True
>>> idade = 70
>>> idade > 60
True
\end{lstlisting}
      \begin{exampleblock}{Ideia-chave}
        O \code{if} não entende ``pressão alta''. Ele só sabe
        olhar para um \code{True} ou um \code{False}.
      \end{exampleblock}
    \end{column}
  \end{columns}
\end{frame}

\begin{frame}[fragile]{\texttt{==} não é \texttt{=}}
  \begin{columns}[T]
    \begin{column}{0.48\textwidth}
      \begin{block}{\code{=} guarda um valor}
\begin{lstlisting}
idade = 70   # agora idade vale 70
\end{lstlisting}
      \end{block}
      \begin{block}{\code{==} faz uma pergunta}
\begin{lstlisting}
idade == 70  # True
\end{lstlisting}
      \end{block}
    \end{column}
    \begin{column}{0.48\textwidth}
      \begin{alertblock}{Erro que o Python recusa}
\begin{lstlisting}
if idade = 70:
    print("idoso")
\end{lstlisting}
        \code{SyntaxError: invalid syntax}
      \end{alertblock}
      \begin{exampleblock}{Como não esquecer}
        Leia em voz alta: \code{=} é ``recebe'', \code{==} é
        ``é igual a?''.
      \end{exampleblock}
    \end{column}
  \end{columns}
\end{frame}

\begin{frame}{\texttt{elif}: ou um, ou outro, nunca dois}
  \begin{center}
  \begin{tikzpicture}[font=\tiny, scale=0.85, transform shape,
    teste/.style={diamond, aspect=3.0, draw=UFSBlue, fill=UFSBlue!12,
      inner sep=0pt, minimum width=3.1cm},
    saida/.style={rectangle, rounded corners, draw=UFSGray,
      fill=UFSLightGray, minimum width=1.7cm, minimum height=0.45cm}]

    \node[teste] (t1) {pressão < 90?};
    \node[teste, below=0.28cm of t1] (t2) {>140 \textbf{e} idoso?};
    \node[teste, below=0.28cm of t2] (t3) {>140 \textbf{ou} idoso?};
    \node[teste, below=0.28cm of t3] (t4) {pressão > 120?};
    \node[saida, below=0.28cm of t4] (t5) {NORMAL};

    \node[saida, right=2.2cm of t1] (s1) {CRITICO};
    \node[saida, right=2.2cm of t2] (s2) {ALTO};
    \node[saida, right=2.2cm of t3] (s3) {MEDIO};
    \node[saida, right=2.2cm of t4] (s4) {BAIXO};

    \foreach \a/\b in {t1/s1, t2/s2, t3/s3, t4/s4}
      \draw[-{Latex[length=1.6mm]}, UFSRed, thick] (\a) -- node[above,font=\tiny]{True} (\b);
    \foreach \a/\b in {t1/t2, t2/t3, t3/t4, t4/t5}
      \draw[-{Latex[length=1.6mm]}, UFSBlue] (\a) -- node[right,font=\tiny]{False} (\b);

    \node[right=0.3cm of s1, font=\tiny, UFSRed, align=left] {sai da função};
  \end{tikzpicture}
  \end{center}

  \vspace{-0.1cm}
  {\small A cascata é uma \textbf{peneira}: o paciente cai no primeiro
  teste que der \code{True} --- os testes de baixo nem chegam a ser
  executados.}
\end{frame}

\begin{frame}[fragile]{\texttt{else}: a rede de segurança}
  \begin{columns}[T]
    \begin{column}{0.52\textwidth}
\begin{lstlisting}
elif pressao > 12.0:
    return "BAIXO"
else:
    return "NORMAL"
\end{lstlisting}
      \vspace{0.2cm}
      O \code{else} não tem condição: ele é ``tudo aquilo que sobrou''.
    \end{column}
    \begin{column}{0.44\textwidth}
      \begin{block}{Por que sempre ter um \code{else}}
        Sem ele, um paciente que não casa com nenhum teste sai da
        função sem classificação --- a função devolve \code{None}
        e o programa imprime \code{Daniela: None}.
      \end{block}
      \begin{alertblock}{Na prova}
        Uma saída \code{None} custa a questão inteira mesmo que toda
        a lógica esteja certa.
      \end{alertblock}
    \end{column}
  \end{columns}
\end{frame}

\begin{frame}{Rastreando o Antonio, linha a linha}
  \code{risco(15.0, 70)}
  \vspace{0.2cm}
  \begin{center}
  \scriptsize
  \begin{tabular}{@{}clll@{}}
    \toprule
    \textbf{\#} & \textbf{Teste avaliado} & \textbf{Resultado} & \textbf{O que acontece} \\
    \midrule
    1 & \code{15.0 < 9.0} & \code{False} & desce para o próximo \\
    2 & \code{15.0 > 14.0 and 70 > 60} & \code{True and True} $\to$ \code{True}
      & \textbf{retorna \code{ALTO}} \\
    3 & \code{15.0 > 14.0 or 70 > 60} & --- & \emph{nunca é executado} \\
    4 & \code{15.0 > 12.0} & --- & \emph{nunca é executado} \\
    5 & \code{else} & --- & \emph{nunca é executado} \\
    \bottomrule
  \end{tabular}
  \end{center}

  \vspace{0.25cm}
  \begin{exampleblock}{A tabela de rastreio é a sua melhor ferramenta}
    Quando o programa dá uma resposta que você não esperava, monte
    esta tabela à mão. Ela mostra exatamente onde a decisão foi tomada.
  \end{exampleblock}
\end{frame}

% ===================================================================
\section{A ordem importa}
% ===================================================================

\begin{frame}[fragile]{Experimento 1: trocar \texttt{elif} por \texttt{if}}
  \begin{columns}[T]
    \begin{column}{0.56\textwidth}
\begin{lstlisting}
def risco(pressao, idade):
    classificacao = "NORMAL"
    if pressao < 9.0:
        classificacao = "CRITICO"
    if pressao > 14.0 and idade > 60:
        classificacao = "ALTO"
    if pressao > 14.0 or idade > 60:
        classificacao = "MEDIO"
    if pressao > 12.0:
        classificacao = "BAIXO"
    return classificacao
\end{lstlisting}
      {\small As mesmas regras, os mesmos números, a mesma ordem.}
    \end{column}
    \begin{column}{0.4\textwidth}
      \begin{alertblock}{Saída}
\begin{lstlisting}[style=saida]
Antonio: BAIXO
Beatriz: BAIXO
Cesar: BAIXO
Daniela: NORMAL
Eduardo: CRITICO
\end{lstlisting}
      \end{alertblock}
      {\small Três pacientes errados. O Antonio, que é \code{ALTO},
      foi para a fila dos casos leves.}
    \end{column}
  \end{columns}
\end{frame}

\begin{frame}{Por que o Antonio virou \texttt{BAIXO}}
  Com \code{if} soltos, \textbf{todos} os testes são executados ---
  não existe peneira. Cada acerto sobrescreve o anterior.

  \vspace{0.25cm}
  \begin{center}
  \scriptsize
  \begin{tabular}{@{}clll@{}}
    \toprule
    \textbf{\#} & \textbf{Teste} & \textbf{Resultado} & \textbf{\code{classificacao} passa a valer} \\
    \midrule
     & --- & --- & \code{"NORMAL"} \\
    1 & \code{15.0 < 9.0} & \code{False} & \code{"NORMAL"} \\
    2 & \code{>14.0 and >60} & \code{True} & \code{"ALTO"} \\
    3 & \code{>14.0 or >60} & \code{True} & \code{"MEDIO"} \\
    4 & \code{15.0 > 12.0} & \code{True} & \code{"BAIXO"} \\
    \bottomrule
  \end{tabular}
  \end{center}

  \vspace{0.25cm}
  \begin{block}{A lição}
    Com \code{if} independentes, quem manda é o \textbf{último} teste
    verdadeiro. Com \code{elif}, quem manda é o \textbf{primeiro}.
  \end{block}
\end{frame}

\begin{frame}[fragile]{Experimento 2: manter o \texttt{elif}, trocar a ordem}
  \begin{columns}[T]
    \begin{column}{0.56\textwidth}
\begin{lstlisting}
def risco(pressao, idade):
    if pressao > 12.0:
        return "BAIXO"
    elif pressao > 14.0 or idade > 60:
        return "MEDIO"
    elif pressao > 14.0 and idade > 60:
        return "ALTO"
    elif pressao < 9.0:
        return "CRITICO"
    else:
        return "NORMAL"
\end{lstlisting}
      {\small A peneira está lá. Só que na ordem errada.}
    \end{column}
    \begin{column}{0.4\textwidth}
      \begin{alertblock}{Saída}
\begin{lstlisting}[style=saida]
Antonio: BAIXO
Beatriz: BAIXO
Cesar: BAIXO
Daniela: NORMAL
Eduardo: CRITICO
\end{lstlisting}
      \end{alertblock}
      \begin{block}{Repare}
        Dois erros \textbf{diferentes} produziram a \textbf{mesma}
        saída errada. Em ambos, o rótulo mais genérico ganhou do
        mais específico.
      \end{block}
    \end{column}
  \end{columns}
\end{frame}

\begin{frame}[shrink=5]{As categorias se contêm umas às outras}
  \begin{center}
  \begin{tikzpicture}[font=\small]
    \draw[fill=UFSYellow!25, draw=UFSGray] (0,0) ellipse (4.4cm and 1.7cm);
    \draw[fill=orange!30, draw=UFSGray]    (0.6,0) ellipse (2.9cm and 1.2cm);
    \draw[fill=UFSRed!35, draw=UFSGray]    (1.4,0) ellipse (1.4cm and 0.7cm);

    \node at (-3.3,0) {\textbf{BAIXO}};
    \node at (-1.1,0) {\textbf{MEDIO}};
    \node at (1.4,0)  {\textbf{ALTO}};

    \node[below, font=\scriptsize] at (0,-1.9)
      {pacientes com \code{pressao > 14.0} e \code{idade > 60} cabem nos três conjuntos};
  \end{tikzpicture}
  \end{center}

  \vspace{0.1cm}
  \begin{block}{A regra geral}
    Quando uma condição é um \textbf{caso particular} de outra, o caso
    particular tem de ser testado \textbf{primeiro}. Senão o teste mais
    largo captura o paciente antes.
  \end{block}
\end{frame}

\begin{frame}{Do mais urgente ao mais geral}
  \begin{center}
  \begin{tikzpicture}[font=\small,
    passo/.style={rectangle, rounded corners, draw=UFSBlue, font=\scriptsize,
      fill=UFSBlue!10, minimum height=0.55cm, text width=7.2cm}]
    \node[passo] (p1) at (0,0)    {\textbf{1. \code{CRITICO}} --- emergência: vale sozinho, ignora tudo};
    \node[passo] (p2) at (0,-0.72) {\textbf{2. \code{ALTO}} --- dois agravantes ao mesmo tempo (\code{and})};
    \node[passo] (p3) at (0,-1.44) {\textbf{3. \code{MEDIO}} --- um agravante qualquer (\code{or})};
    \node[passo] (p4) at (0,-2.16) {\textbf{4. \code{BAIXO}} --- só um sinal leve};
    \node[passo] (p5) at (0,-2.88) {\textbf{5. \code{NORMAL}} --- o que sobrou (\code{else})};

    \draw[-{Latex[length=2.5mm]}, UFSRed, very thick] (-4.5,0.25) -- (-4.5,-3.15);
    \node[rotate=90, UFSRed, font=\tiny] at (-4.8,-1.45)
      {mais específico $\to$ mais geral};
  \end{tikzpicture}
  \end{center}

  \vspace{0.05cm}
  \begin{exampleblock}{Como descobrir a ordem sozinho}
    Pergunte de cada par: ``existe paciente que satisfaz as duas?''
    Se existir, a que descreve \emph{menos} pacientes vem antes.
  \end{exampleblock}
\end{frame}

\begin{frame}[fragile,shrink=5]{\texttt{and}, \texttt{or}, \texttt{not}}
  \begin{columns}[T]
    \begin{column}{0.5\textwidth}
      \begin{center}
      \small
      \begin{tabular}{@{}llll@{}}
        \toprule
        \textbf{A} & \textbf{B} & \code{A and B} & \code{A or B} \\
        \midrule
        \code{True}  & \code{True}  & \code{True}  & \code{True} \\
        \code{True}  & \code{False} & \code{False} & \code{True} \\
        \code{False} & \code{True}  & \code{False} & \code{True} \\
        \code{False} & \code{False} & \code{False} & \code{False} \\
        \bottomrule
      \end{tabular}
      \end{center}
      \vspace{0.2cm}
      {\small \code{not True} $\to$ \code{False} \hspace{0.4cm}
      \code{not False} $\to$ \code{True}}
    \end{column}
    \begin{column}{0.46\textwidth}
      \begin{block}{No enunciado da triagem}
        \code{ALTO}: pressão elevada \textbf{e} paciente idoso
        $\to$ \code{and} \\[0.25cm]
        \code{MEDIO}: pressão elevada \textbf{ou} paciente idoso
        $\to$ \code{or}
      \end{block}
      \begin{alertblock}{Armadilha frequente}
\begin{lstlisting}
if pressao > 14.0 and > 60:
\end{lstlisting}
        Não existe. Cada lado do \code{and} precisa ser uma condição
        completa: \code{pressao > 14.0 and idade > 60}.
      \end{alertblock}
    \end{column}
  \end{columns}
\end{frame}

\begin{frame}[fragile]{Cuidado com as fronteiras}
  ``Acima de 140'' é \code{> 14.0}. Não é \code{>= 14.0}.
  Um caractere muda a resposta.

  \vspace{0.25cm}
  \begin{center}
  \scriptsize
  \begin{tabular}{@{}llll@{}}
    \toprule
    \textbf{Pressão} & \textbf{Idade} & \textbf{Classificação} & \textbf{Por quê} \\
    \midrule
    \code{14.0} & \code{70} & \code{MEDIO}  & 140 \emph{não} é acima de 140; mas 70 é acima de 60 \\
    \code{14.0} & \code{30} & \code{BAIXO}  & nenhum agravante; 140 está acima de 120 \\
    \code{12.0} & \code{20} & \code{NORMAL} & 120 \emph{não} é acima de 120 \\
    \code{9.0}  & \code{45} & \code{NORMAL} & 90 \emph{não} é inferior a 90 \\
    \code{8.99} & \code{45} & \code{CRITICO}& agora sim, abaixo de 90 \\
    \bottomrule
  \end{tabular}
  \end{center}

  \vspace{0.25cm}
  \begin{block}{Hábito de quem não perde ponto}
    Antes de entregar, teste os valores que ficam \emph{exatamente} em
    cima do limite. É lá que o erro se esconde.
  \end{block}
\end{frame}

% ===================================================================
\section{Sua vez}
% ===================================================================

\begin{frame}[fragile,shrink=8]{Modifique 1 --- conserte a versão quebrada}
  \begin{columns}[T]
    \begin{column}{0.55\textwidth}
\begin{lstlisting}[basicstyle=\ttfamily\tiny]
def risco(pressao, idade):
    if pressao > 12.0:
        return "BAIXO"
    elif pressao > 14.0 or idade > 60:
        return "MEDIO"
    elif pressao > 14.0 and idade > 60:
        return "ALTO"
    elif pressao < 9.0:
        return "CRITICO"
    else:
        return "NORMAL"
\end{lstlisting}
    \end{column}
    \begin{column}{0.42\textwidth}
      \begin{block}{Sua tarefa}
        {\small Reordene os testes --- sem mudar nenhuma condição e
        sem acrescentar linha --- para que a saída seja:}
\begin{lstlisting}[style=saida, basicstyle=\ttfamily\tiny]
Antonio: ALTO
Beatriz: MEDIO
Cesar: MEDIO
Daniela: NORMAL
Eduardo: CRITICO
\end{lstlisting}
      \end{block}
      \begin{exampleblock}{Antes de rodar}
        {\small Diga em voz alta qual linha tem de subir e por quê.}
      \end{exampleblock}
    \end{column}
  \end{columns}
\end{frame}

\begin{frame}[shrink=8]{Modifique 2 --- o limite exato}
  \begin{block}{Cenário}
    {\small O protocolo mudou: pressão \textbf{de} 140 mmHg
    \textbf{para cima} (e não apenas acima de 140) já conta como
    pressão elevada. A idade continua ``acima de 60''.}
  \end{block}

  \vspace{0.15cm}
  \begin{block}{Sua tarefa}
    \small
    \begin{enumerate}
      \item Faça a alteração no código.
      \item \textbf{Antes de rodar}, preveja: \emph{algum} dos cinco
            pacientes muda de classificação? Quais?
      \item Rode e confira. Se errou a previsão, monte a tabela de
            rastreio desse paciente.
    \end{enumerate}
  \end{block}

  \vspace{0.15cm}
  \begin{exampleblock}{Para pensar}
    {\small Quantos caracteres você mudou? E quantos pacientes isso
    afetaria num hospital com mil atendimentos por dia?}
  \end{exampleblock}
\end{frame}

\begin{frame}[shrink=8]{Modifique 3 --- uma categoria nova}
  \begin{block}{Cenário}
    {\small A pediatria pediu uma sexta classificação:
    \code{PEDIATRICO} --- paciente com \textbf{menos de 12 anos} e
    pressão acima de 120 mmHg. Quem está \code{CRITICO} continua
    \code{CRITICO}, em qualquer idade.}
  \end{block}

  \vspace{0.15cm}
  \begin{block}{Sua tarefa}
    \small
    \begin{enumerate}
      \item Em que posição da cascata o novo \code{elif} entra?
            Justifique antes de escrever.
      \item Implemente e teste com \code{(13.0, 8)}, \code{(11.0, 8)},
            \code{(8.0, 5)} e \code{(13.0, 40)}.
      \item Confira que os cinco pacientes originais não mudaram.
    \end{enumerate}
  \end{block}

  \vspace{0.15cm}
  \begin{alertblock}{Critério de qualidade}
    {\small Uma alteração bem feita não quebra o que já funcionava.}
  \end{alertblock}
\end{frame}

\begin{frame}[fragile,shrink=6]{Crie --- o programa completo da Q1}
  \begin{columns}[T]
    \begin{column}{0.44\textwidth}
      \begin{block}{O que o programa faz}
        {\small Lê \code{N}. Para cada paciente, lê nome, pressão e
        idade (uma por linha). Imprime \code{<nome>: <classificacao>}.}
      \end{block}
      \begin{alertblock}{Formato exato}
        {\small Dois-pontos, um espaço, classificação em maiúsculas.
        Nem um espaço a mais.}
      \end{alertblock}
    \end{column}
    \begin{column}{0.52\textwidth}
\begin{lstlisting}
def risco(pressao, idade):
    # sua cascata de if / elif / else

n = int(input())
for _ in range(n):
    nome = input()
    pressao = float(input())
    idade = int(input())
    print(nome + ": " + risco(pressao, idade))
\end{lstlisting}
      \begin{exampleblock}{Atenção aos tipos}
        {\small \code{input()} devolve texto. Pressão é \code{float},
        idade é \code{int}. Sem a conversão, comparar texto com número
        levanta \code{TypeError}.}
      \end{exampleblock}
    \end{column}
  \end{columns}
\end{frame}

\begin{frame}[shrink=4]{Crie --- leve a ideia para outro domínio}
  \begin{block}{Desafio}
    \small
    Escreva \code{classifica\_chuva(mm, horas)} que classifica a
    precipitação de uma cidade:
    \begin{itemize}
      \item \code{ENCHENTE}: mais de 100 mm em menos de 6 horas
      \item \code{ALERTA}: mais de 100 mm (em qualquer tempo)
            \textbf{ou} mais de 50 mm em menos de 6 horas
      \item \code{ATENCAO}: mais de 50 mm
      \item \code{NORMAL}: os demais casos
    \end{itemize}
  \end{block}

  \vspace{0.15cm}
  \begin{exampleblock}{O que se pede de verdade}
    {\small Antes de programar, escreva no papel a \textbf{ordem} dos
    testes e aponte qual condição é caso particular de qual. O código
    é a parte fácil; a ordem é a parte que vale a nota.}
  \end{exampleblock}
\end{frame}

\begin{frame}{O que costuma custar ponto}
  \begin{center}
  \scriptsize
  \begin{tabular}{@{}p{5.6cm}p{6.4cm}@{}}
    \toprule
    \textbf{Erro} & \textbf{Como evitar} \\
    \midrule
    Testar do mais geral para o mais específico
      & Comece pela emergência, termine no \code{else} \\
    Usar \code{if} solto onde cabia \code{elif}
      & Se as respostas são exclusivas, use \code{elif} \\
    Trocar \code{>} por \code{>=} nos limites
      & Teste os valores exatamente em cima do limite \\
    Esquecer o \code{else}
      & Sem ele a função devolve \code{None} \\
    Esquecer \code{int()} / \code{float()} no \code{input()}
      & Comparar texto com número dá \code{TypeError} \\
    Alterar o formato da saída
      & Copie o formato do enunciado, caractere a caractere \\
    \bottomrule
  \end{tabular}
  \end{center}
\end{frame}

\begin{frame}{Para levar para casa}
  \begin{block}{Os quatro fatos}
    \begin{enumerate}
      \item Uma condição é uma expressão que vale \code{True} ou \code{False}.
      \item \code{elif} garante que \textbf{exatamente uma} alternativa é
            escolhida --- a primeira verdadeira.
      \item Quando as descrições se sobrepõem, a \textbf{ordem} dos testes
            é que define a resposta: do mais específico ao mais geral.
      \item Os erros moram nas fronteiras. Teste os limites.
    \end{enumerate}
  \end{block}

  \vspace{0.2cm}
  \begin{exampleblock}{Antes da próxima aula}
    Termine no notebook: a triagem com a categoria \code{PEDIATRICO} e
    a função \code{classifica\_chuva}. Traga a sua tabela de rastreio de
    um caso que você errou --- vamos discutir em sala.
  \end{exampleblock}
\end{frame}

\section*{Referências}
\begin{frame}[allowframebreaks]{Referências}
  \scriptsize
  \nocite{*}
  \bibliographystyle{plain}
  \bibliography{referencias}
\end{frame}

\end{document}
